\subsection{Selection of Decision Scenarios}
\label{sec:scenarios}

We evaluate eight sequential decision making scenarios that share the
same Bayesian \(\mathrm{BASELINE}\) procedure from
Section~\ref{sec:baseline-design} but span a wide range of decision
structures. They differ in their hidden parameter space \(\Theta_S\),
observation model \(P_S(o_s\mid a_s,\theta)\), and utility function
\(U_S(a,\theta)\), and they range from inference about a property of a
nonstrategic environment to inference about the strategy of a
history-dependent opponent. We first define their common generative
process and design requirements.

\textbf{\textit{Common generative template.}}
Every scenario is instantiated by an \emph{environment generator} \(G_S\).
For a scenario \(S\) with finite hidden parameter space \(\Theta_S\), each
rollout is generated in three steps. First, \(G_S\) samples a hidden
parameter \(\theta\sim\mathrm{Unif}(\Theta_S)\) once and holds it fixed for
the full horizon \(T=10\). Second, at each step \(s\), a fixed behavior policy
selects one of the two available actions \(a_s\) with equal probability,
independently of the evaluated model. Third, the environment returns an
outcome \(o_s\sim P_S(o_s\mid a_s,\theta)\). The history available at step
\(t\) is \(h_t=((a_1,o_1),\ldots,(a_{t-1},o_{t-1}))\). Because histories
are generated by the fixed behavior policy rather than by the evaluated
model, all models receive the same evidence trajectory for a given seed.

\textbf{\textit{Design requirements.}}
We design each scenario to satisfy the following requirements:

\begin{itemize}
    \item[(C1)] A finite hidden parameter space
    \(\Theta_S\) with a known uniform prior.

    \item[(C2)] Sequential observations that are informative about
    \(\theta\), allowing the posterior \(P(\theta\mid h_t)\) to update as
    evidence accumulates.

    \item[(C3)] A binary action space.

    \item[(C4)] An analytically computable rational decision rule from the
    posterior belief and scenario utility, with indifference cases
    explicitly identified.

    \item[(C5)] A posterior that is decision relevant, in the sense that
    some belief state makes each of the two actions optimal.
\end{itemize}

These properties allow both model behavior and internal representations to
be evaluated against the same \(\mathrm{BASELINE}\) while varying
\(\Theta_S\), \(P_S(o_s\mid a_s,\theta)\), and \(U_S(a,\theta)\).

\textbf{\textit{Bayesian inference family.}}
In this family, the hidden parameter is a property of a nonstrategic
environment and \(P_S(o_s\mid a_s,\theta)\) is a fixed stochastic rule.
We separate these scenarios into pure inference tasks, whose decision
boundary is a posterior probability of \(0.5\), and expected utility
threshold tasks, whose decision boundary depends on the posterior mean of
a latent rate \(q\). The latter tests whether a model can recover a
nontrivial expected utility threshold rather than only perform binary
hypothesis testing.

\begin{itemize}
    \item \textit{Bandit.}
    A two-armed Bernoulli bandit in which the identity of the better arm is
    hidden and inferred from action-dependent rewards.

    \item \textit{Monty Hall.}
    An altered Monty Hall problem with a hidden biased door placement rule.
    The model must update away from the standard switching probability when
    the observed history supports a different rule.

    \item \textit{Trust Game.}
    A repeated investor-side trust game in which the trustee has a hidden
    repayment rate \(q\), and investing is optimal only above a nontrivial
    expected utility threshold.

    \item \textit{Machine.}
    A nonsocial counterpart to the Trust Game with matched hidden types,
    utilities, and outcome trajectories, isolating social framing while
    preserving the underlying inference problem.
\end{itemize}

\textbf{\textit{Game theoretic family.}}
In this family, \(U_S(a,\theta)\) is induced by a payoff matrix over the
joint action profile, and the hidden parameter \(\theta\) is a stationary
property of the opponent: either its strategy type or its action rate.
Since the opponent is drawn from a known prior and inferred through
Bayesian updating, these tasks are Bayesian decision problems against fixed
opponent types rather than equilibrium computations. The opponent strategy
remains stationary throughout each rollout. Three of the four scenarios
include a reactive type whose next action depends on the agent's previous
action, so the decision relevant quantity is the posterior predictive over
the opponent's next move rather than the posterior over types alone. The
three also place the decision boundary at a different posterior threshold,
so a model cannot succeed across them by applying one fixed cutoff.

\begin{itemize}
    \item \textit{Stag Hunt.}
    A repeated coordination game against a hidden stationary opponent type,
    including a reactive tit-for-tat strategy. Coordinating on stag is
    optimal only when the posterior predictive probability of a stag move
    exceeds \(2/3\).

    \item \textit{Bach or Stravinsky.}
    A repeated coordination game with conflicting preferences over the two
    coordination points, against the same hidden type structure including
    tit-for-tat. Because the two coordination outcomes pay the model
    differently, the decision boundary is an asymmetric threshold of
    \(1/3\) rather than \(0.5\).

    \item \textit{Matching Pennies.}
    A strictly competitive matching game whose hidden types include an
    anti-imitative counter strategy that plays the opposite of the agent's
    previous action. No action is dominant, so the optimal action tracks
    the posterior predictive over the opponent's next move around a
    threshold of \(0.5\).

    \item \textit{Chicken.}
    A repeated anti-coordination game against an opponent with a hidden
    stationary yield rate.
\end{itemize}

Table~\ref{tab:decision-scenarios} summarizes, for each scenario, the
hidden parameter and observation model, binary action space, utility or
payoff structure, and rational decision rule.
